\subsection{The open projection to unmeasured spaces}\label{subsec:measures-open-projection}
\paraid{p-fef9f75c2978}
We next restrict the measured space
$\MeasuredSpaces$
to invariant full-support measures
and prove that forgetting the measure is an open map.
We use uniform control of integrals to prove that the full-support condition
is a countable intersection of open conditions.
By taking limits of isometries,
we prove that the invariance condition is also
a countable intersection of open conditions and that the averaged measures
used to lift a convergent sequence of compact spaces converge weakly.
The required uniform estimates are in
\Cref{lem:uniform-weak-integrals,lem:parameter-integrals}.

\paraid{p-573911841fb0}
In the next proposition,
we extract a convergent subsequence of graphs of isometries
when the compact spaces vary.

\begin{proposition}\label{lem:isometry-limits}
\paraid{p-7c4e086d9e70}
Let
$\BaseCarrier _\SequenceIndex ,\BaseCarrier \subset \AmbientCarrier $
be nonempty compact subsets of a compact metric space
$\MetricSpace{\AmbientCarrier }{\AmbientMetric }$
with
$\HausdorffDistance{\AmbientMetric }(\BaseCarrier _\SequenceIndex ,\BaseCarrier )\to0$.
Equip each subset of
$\AmbientCarrier$
with the restricted ambient metric, and equip the product
$\AmbientCarrier\times\AmbientCarrier$
with the maximum product metric.
For any isometries
$\Isometry _\SequenceIndex \colon \BaseCarrier _\SequenceIndex \to \BaseCarrier _\SequenceIndex $,
a subsequence of their graphs converges in
$\CompactHyperspace(\AmbientCarrier \times \AmbientCarrier )$
to the graph of an isometry
$\Isometry \colon \BaseCarrier \to \BaseCarrier $.
Along that subsequence,
let
$\GraphProbability_\SequenceIndex\in\ProbabilityMeasures(\BaseCarrier_\SequenceIndex)$
and
$\GraphProbability\in\ProbabilityMeasures(\BaseCarrier)$
satisfy
$\GraphProbability_\SequenceIndex\to\GraphProbability$
weakly as probabilities on
$\AmbientCarrier$.
Then
\[
 (\Isometry_\SequenceIndex)\Pushforward\GraphProbability_\SequenceIndex
 \longrightarrow\Isometry\Pushforward\GraphProbability
 \quad\text{weakly in }\ProbabilityMeasures(\AmbientCarrier).
\]
\end{proposition}

\begin{proof}
\paraid{p-b6a94ad3a354}
We first identify the limit of the graphs and then the limit of the
probabilities supported on them.
For all
$\BasePoint_0,\ComparisonPoint_0,\BasePoint_1,\ComparisonPoint_1\in\AmbientCarrier$,
write
\[
 \AmbientMetric_\times((\BasePoint_0,\ComparisonPoint_0),(\BasePoint_1,\ComparisonPoint_1))
 =\max\{\AmbientMetric(\BasePoint_0,\BasePoint_1),
          \AmbientMetric(\ComparisonPoint_0,\ComparisonPoint_1)\}.
\]
For each
$\SequenceIndex$,
put
\[
 \Correspondence_\SequenceIndex
 =\{(\BasePoint,\Isometry_\SequenceIndex(\BasePoint))\mid
       \BasePoint\in\BaseCarrier_\SequenceIndex\}.
\]
By compactness of
$\CompactHyperspace(\AmbientCarrier\times\AmbientCarrier)$,
choose a subsequence whose graphs converge to a nonempty compact set
$\Correspondence\subset\AmbientCarrier\times\AmbientCarrier$.
For the remainder of this proof,
reindex the selected spaces,
isometries,
probabilities,
and graphs by
$\SequenceIndex\in\NonnegativeIntegers$,
retaining the notation
$\BaseCarrier_\SequenceIndex$,
$\Isometry_\SequenceIndex$,
$\GraphProbability_\SequenceIndex$,
and
$\Correspondence_\SequenceIndex$.
Then
\begin{equation}\label{eq:isometry-graph-convergence}
 \HausdorffDistance{\AmbientMetric_\times}
 (\Correspondence_\SequenceIndex,\Correspondence)\longrightarrow0.
\end{equation}

\textbf{Step 1.
The limiting graph.}
\paraid{p-7e16f5aad565}
Let
$\CoordinateProjection_1,\CoordinateProjection_2\colon
 \AmbientCarrier\times\AmbientCarrier\to\AmbientCarrier$
be the coordinate projections.
They are
$1$-Lipschitz
for
$\AmbientMetric_\times$,
and
$\CoordinateProjection_i(\Correspondence_\SequenceIndex)=\BaseCarrier_\SequenceIndex$
for
$i=1,2$,
because each
$\Isometry_\SequenceIndex$
is surjective.
For each
$i=1,2$,
we apply \Cref{lem:hausdorff-lipschitz-images} to
$\CoordinateProjection_i$
with
$L=1$.
Together with the triangle inequality,
this yields
\[
 \begin{aligned}
 \HausdorffDistance{\AmbientMetric}(\CoordinateProjection_i(\Correspondence),\BaseCarrier)
 &\leq\HausdorffDistance{\AmbientMetric}
       (\CoordinateProjection_i(\Correspondence),\CoordinateProjection_i(\Correspondence_\SequenceIndex))
       +\HausdorffDistance{\AmbientMetric}(\BaseCarrier_\SequenceIndex,\BaseCarrier)\\
 &\leq\HausdorffDistance{\AmbientMetric_\times}(\Correspondence,\Correspondence_\SequenceIndex)
       +\HausdorffDistance{\AmbientMetric}(\BaseCarrier_\SequenceIndex,\BaseCarrier)
       \longrightarrow0.
 \end{aligned}
\]
Thus
$\CoordinateProjection_1(\Correspondence)=\CoordinateProjection_2(\Correspondence)=\BaseCarrier$.
In particular,
$\Correspondence\subset\BaseCarrier\times\BaseCarrier$.

\paraid{p-49be4c2dec12}
Fix
$(\BasePoint_0,\ComparisonPoint_0),(\BasePoint_1,\ComparisonPoint_1)\in\Correspondence$.
For each
$i\in\{0,1\}$
and each
$\SequenceIndex\in\NonnegativeIntegers$,
choose a pair
$(\BasePoint_{i,\SequenceIndex},\ComparisonPoint_{i,\SequenceIndex})\in\Correspondence_\SequenceIndex$
whose distance from
$(\BasePoint_i,\ComparisonPoint_i)$
in the metric
$\AmbientMetric_\times$
is minimal.
Such a pair exists because
$\Correspondence_\SequenceIndex$
is nonempty and compact.
By \eqref{eq:isometry-graph-convergence},
\[
 \begin{aligned}
 \max\{\AmbientMetric(\BasePoint_{i,\SequenceIndex},\BasePoint_i),
        \AmbientMetric(\ComparisonPoint_{i,\SequenceIndex},\ComparisonPoint_i)\}
 &=\dist_{\AmbientMetric_\times}
   ((\BasePoint_i,\ComparisonPoint_i),\Correspondence_\SequenceIndex)\\
 &\leq\HausdorffDistance{\AmbientMetric_\times}
   (\Correspondence_\SequenceIndex,\Correspondence)\longrightarrow0.
 \end{aligned}
\]
Thus
$\BasePoint_{i,\SequenceIndex}\to\BasePoint_i$
and
$\ComparisonPoint_{i,\SequenceIndex}\to\ComparisonPoint_i$.
Since
$\ComparisonPoint_{i,\SequenceIndex}=\Isometry_\SequenceIndex(\BasePoint_{i,\SequenceIndex})$,
we obtain
\begin{equation}\label{eq:limit-graph-distance}
 \begin{aligned}
 \AmbientMetric(\BasePoint_0,\BasePoint_1)
 &=\lim_\SequenceIndex\AmbientMetric(\BasePoint_{0,\SequenceIndex},\BasePoint_{1,\SequenceIndex})\\
 &=\lim_\SequenceIndex\AmbientMetric(
      \Isometry_\SequenceIndex(\BasePoint_{0,\SequenceIndex}),
      \Isometry_\SequenceIndex(\BasePoint_{1,\SequenceIndex}))\\
 &=\lim_\SequenceIndex\AmbientMetric(\ComparisonPoint_{0,\SequenceIndex},\ComparisonPoint_{1,\SequenceIndex})
 =\AmbientMetric(\ComparisonPoint_0,\ComparisonPoint_1).
 \end{aligned}
\end{equation}
In particular,
if
 $(\BasePoint,\ComparisonPoint_0),(\BasePoint,\ComparisonPoint_1)\in\Correspondence$,
 then
 $\AmbientMetric(\ComparisonPoint_0,\ComparisonPoint_1)=0$.
 Hence
 $\ComparisonPoint_0=\ComparisonPoint_1$.
Since the first projection is surjective and each first coordinate has a unique
second coordinate,
we define the map
$\Isometry\colon\BaseCarrier\to\BaseCarrier$
by
\[
 \Correspondence=\{(\BasePoint,\Isometry(\BasePoint))\mid\BasePoint\in\BaseCarrier\}.
\]
Since the second projection is surjective,
we have
$\Isometry(\BaseCarrier)=\BaseCarrier$,
and for every
$\BasePoint_0,\BasePoint_1\in\BaseCarrier$,
\eqref{eq:limit-graph-distance} implies
\[
 \AmbientMetric(\Isometry(\BasePoint_0),\Isometry(\BasePoint_1))
 =\AmbientMetric(\BasePoint_0,\BasePoint_1).
\]
Hence
$\Isometry$
is an isometry.

\textbf{Step 2.
Probabilities on the graphs.}
\paraid{p-338e44ba371c}
For every
$\SequenceIndex\in\NonnegativeIntegers$,
define the map
\[
 \GraphMap_\SequenceIndex\colon\BaseCarrier_\SequenceIndex\to\AmbientCarrier\times\AmbientCarrier
\]
by
\[
 \GraphMap_\SequenceIndex(\BasePoint)=(\BasePoint,\Isometry_\SequenceIndex(\BasePoint)).
\]
We also define the map
\[
 \GraphMap\colon\BaseCarrier\to\AmbientCarrier\times\AmbientCarrier
\]
by
\[
 \GraphMap(\BasePoint)=(\BasePoint,\Isometry(\BasePoint)).
\]
Set
\[
 \GraphLaw_\SequenceIndex=(\GraphMap_\SequenceIndex)\Pushforward\GraphProbability_\SequenceIndex.
\]
By compactness of
$\ProbabilityMeasures(\AmbientCarrier\times\AmbientCarrier)$,
every subsequence of
$\GraphLaw_\SequenceIndex$
has a further weakly convergent subsequence.
Write the indices of any such subsequence as
$n_k$
and denote its limit by
$\GraphLaw\in\ProbabilityMeasures(\AmbientCarrier\times\AmbientCarrier)$.
Then
\[
 \GraphLaw_{n_k}\longrightarrow\GraphLaw
 \quad\text{weakly in }\ProbabilityMeasures(\AmbientCarrier\times\AmbientCarrier).
\]
Since
$\GraphLaw_{n_k}(\Correspondence_{n_k})=1$,
Equation \eqref{eq:isometry-graph-convergence},
\Cref{lem:limit-support},
and continuity of the first projection
$\CoordinateProjection_1$
imply
\begin{equation}\label{eq:limit-graph-law-support}
 \GraphLaw(\Correspondence)=1,
\end{equation}
and
\[
 (\CoordinateProjection_1)\Pushforward\GraphLaw
 =\lim_{k\to\infty}(\CoordinateProjection_1)\Pushforward\GraphLaw_{n_k}
 =\lim_{k\to\infty}\GraphProbability_{n_k}
 =\GraphProbability.
\]
On
$\Correspondence$
we have
\begin{equation}\label{eq:graph-projection-inverse}
 \GraphMap\circ\CoordinateProjection_1=\operatorname{id}_\Correspondence.
\end{equation}
Regard
$\GraphLaw$
as a probability on
$\Correspondence$.
By \eqref{eq:graph-projection-inverse},
we obtain
\[
 \GraphLaw
 =(\GraphMap\circ\CoordinateProjection_1)\Pushforward\GraphLaw
 =\GraphMap\Pushforward\bigl((\CoordinateProjection_1)\Pushforward\GraphLaw\bigr)
 =\GraphMap\Pushforward\GraphProbability.
\]
Every weakly convergent subsequence has this same limit.
Compactness of the probability space
$\ProbabilityMeasures(\AmbientCarrier\times\AmbientCarrier)$
consequently implies
\[
 \GraphLaw_\SequenceIndex\longrightarrow\GraphMap\Pushforward\GraphProbability
 \quad\text{weakly in }\ProbabilityMeasures(\AmbientCarrier\times\AmbientCarrier).
\]
Applying the continuous second projection
$\CoordinateProjection_2$
yields
\[
 (\Isometry_\SequenceIndex)\Pushforward\GraphProbability_\SequenceIndex
 =(\CoordinateProjection_2)\Pushforward\GraphLaw_\SequenceIndex
 \longrightarrow(\CoordinateProjection_2)\Pushforward(\GraphMap\Pushforward\GraphProbability)
 =\Isometry\Pushforward\GraphProbability.
\]
This is the required weak convergence on
$\AmbientCarrier$.
\end{proof}

\begin{proposition}\label{lem:full-support-gdelta}
\paraid{p-380dfd13f917}
The subset
\[
 \mathscr S_{\mathrm{full}}
 =\{\MeasuredSpace{\BaseCarrier}{\MetricSymbol}{\FirstMeasure}\in\MeasuredSpaces
       \mid\supp(\FirstMeasure)=\BaseCarrier\}
\]
is a
$\CountableIntersectionType$
subset of
$\MeasuredSpaces$.
\end{proposition}

\begin{proof}
\paraid{p-0aa6b81d8d5c}
For
$\MetricSpace{\BaseCarrier}{\MetricSymbol}\in\GHSpace$,
$\FirstMeasure\in\ProbabilityMeasures(\BaseCarrier)$,
and
$\Radius>0$,
define the ball mass function at each
$\BasePoint\in\BaseCarrier$
\[
 \BallMassFunction_\Radius\colon\BaseCarrier\to[0,1]
\]
by
\begin{equation}\label{eq:ball-mass-functions-1}
 \BallMassFunction_\Radius(\BasePoint)
 =\int_\BaseCarrier\BumpFunction{\BaseCarrier}{\MetricSymbol}{\BasePoint}{\Radius}
       \,\IntegrationDifferential\FirstMeasure,
\end{equation}
and define
\[
 \MinimumBallMass_\Radius\colon\MeasuredSpaces\to[0,1]
\]
by
\begin{equation}\label{eq:ball-mass-functions-2}
 \MinimumBallMass_\Radius\MeasuredSpace{\BaseCarrier}{\MetricSymbol}{\FirstMeasure}
 =\min_{\BasePoint\in\BaseCarrier}\BallMassFunction_\Radius(\BasePoint).
\end{equation}
By \eqref{eq:metric-bump-lipschitz},
the function
$\BallMassFunction_\Radius$
is
$1/\Radius$-Lipschitz,
so its minimum exists.
To confirm the continuity of
$\MinimumBallMass_\Radius$,
let
$\MeasuredSpace{\BaseCarrier_\SequenceIndex}{\MetricSymbol_\SequenceIndex}{\FirstMeasure_\SequenceIndex}$
converge to
$\MeasuredSpace{\BaseCarrier}{\MetricSymbol}{\FirstMeasure}$
in
$\MeasuredSpaces$.
Use \Cref{lem:common-measured-embedding} to realize the carriers in
$\MetricSpace{\AmbientCarrier}{\AmbientMetric}$
with Hausdorff convergence and
$\FirstMeasure_\SequenceIndex\to\FirstMeasure$
weakly on
$\AmbientCarrier$.
Define ambient extensions of the ball mass functions
by integrating the ambient bump functions against the measures supported on the respective carriers.
For every
$\IntegrationPoint\in\AmbientCarrier$,
define
\begin{equation}\label{eq:ambient-ball-mass-functions-1}
 \BallMassFunction_{\Radius,\SequenceIndex}(\IntegrationPoint)
 =\int_\AmbientCarrier
    \BumpFunction{\AmbientCarrier}{\AmbientMetric}{\IntegrationPoint}{\Radius}
    \,\IntegrationDifferential\FirstMeasure_\SequenceIndex,
\end{equation}
and
\begin{equation}\label{eq:ambient-ball-mass-functions-2}
 \tilde{\BallMassFunction}_\Radius(\IntegrationPoint)
 =\int_\AmbientCarrier
    \BumpFunction{\AmbientCarrier}{\AmbientMetric}{\IntegrationPoint}{\Radius}
    \,\IntegrationDifferential\FirstMeasure.
\end{equation}
We have
$\tilde{\BallMassFunction}_\Radius|_\BaseCarrier=\BallMassFunction_\Radius$,
and
$\BallMassFunction_{\Radius,\SequenceIndex}|_{\BaseCarrier_\SequenceIndex}$
is the ball mass function in \eqref{eq:ball-mass-functions-1} for
$\MeasuredSpace{\BaseCarrier_\SequenceIndex}{\MetricSymbol_\SequenceIndex}{\FirstMeasure_\SequenceIndex}$.
By \Cref{lem:parameter-integrals} and \eqref{eq:metric-bump-lipschitz},
\begin{equation}\label{eq:uniform-ball-mass-convergence}
 \norm{\BallMassFunction_{\Radius,\SequenceIndex}-\tilde{\BallMassFunction}_\Radius}_\infty
 \longrightarrow0.
\end{equation}
Using \eqref{eq:uniform-ball-mass-convergence} and the
$1/\Radius$-Lipschitz bound,
we obtain
\begin{equation}\label{eq:minimum-ball-mass-convergence}
 |\min_{\BasePoint\in\BaseCarrier_\SequenceIndex}\BallMassFunction_{\Radius,\SequenceIndex}(\BasePoint)
      -\min_{\BasePoint\in\BaseCarrier}\tilde{\BallMassFunction}_\Radius(\BasePoint)|
 \leq\norm{\BallMassFunction_{\Radius,\SequenceIndex}-\tilde{\BallMassFunction}_\Radius}_\infty
      +\frac{\HausdorffDistance{\AmbientMetric}(\BaseCarrier_\SequenceIndex,\BaseCarrier)}{\Radius}
 \longrightarrow0.
\end{equation}
Thus
$\MinimumBallMass _\Radius $
is continuous on
$\MeasuredSpaces$.

\paraid{p-871abef1344e}
If
$\supp(\FirstMeasure )=\BaseCarrier $,
then \eqref{eq:metric-bump-integral} implies that for every
$\BasePoint \in \BaseCarrier $,
\[
 \BallMassFunction _\Radius (\BasePoint )\geq\tfrac12\FirstMeasure (\MetricBall(\BasePoint ,\Radius /2;\MetricSymbol))>0.
\]
Continuity and compactness imply
$\MinimumBallMass _\Radius >0$.
Conversely,
if
$\BasePoint \notin\supp(\FirstMeasure )$,
choose $r>0$ such that
$\FirstMeasure(\MetricBall(\BasePoint,r;\MetricSymbol))=0$.
Choose $\CoordinateIndex\in\NonnegativeIntegers$ with $2^{-\CoordinateIndex}<r$.
Then $\FirstMeasure(\MetricBall(\BasePoint,2^{-\CoordinateIndex};\MetricSymbol))=0$.
Then
$\BallMassFunction _{2^{-\CoordinateIndex} }(\BasePoint )=0$.
We have proved that
$\supp(\FirstMeasure)=\BaseCarrier$
if and only if for every
$\CoordinateIndex\in\NonnegativeIntegers$
we have
\[
 \MinimumBallMass_{2^{-\CoordinateIndex}}
 \MeasuredSpace{\BaseCarrier}{\MetricSymbol}{\FirstMeasure}>0.
\]
For each
$\CoordinateIndex\in\NonnegativeIntegers$,
define the subset
\[
 \PositiveBallMassSet _\CoordinateIndex =\{\FiberPoint \in\MeasuredSpaces\mid \MinimumBallMass _{2^{-\CoordinateIndex} }(\FiberPoint )>0\}.
\]
Each
$\PositiveBallMassSet_\CoordinateIndex$
is open by continuity of
$\MinimumBallMass_{2^{-\CoordinateIndex}}$.

\paraid{p-39a712fdc965}
Consequently,
\[
 \mathscr S_{\mathrm{full}}
 =\bigcap_{\CoordinateIndex\in\NonnegativeIntegers}\PositiveBallMassSet_\CoordinateIndex.
\]
Thus
$\mathscr S_{\mathrm{full}}$
is a countable intersection of open subsets of
$\MeasuredSpaces$,
as required.
\end{proof}

\begin{proposition}\label{lem:invariant-measures-gdelta}
\paraid{p-fc58c5a6a993}
Let
$\mathscr S_{\mathrm{inv}}\subset\MeasuredSpaces$
consist of the classes
$\MeasuredSpace{\BaseCarrier}{\MetricSymbol}{\FirstMeasure}$
whose measures are invariant under every isometry of
$\MetricSpace{\BaseCarrier}{\MetricSymbol}$.
Then
$\mathscr S_{\mathrm{inv}}$
is a
$\CountableIntersectionType$
subset of
$\MeasuredSpaces$.
\end{proposition}

\begin{proof}
\paraid{p-9a02df523b4a}
Fix
$\CoordinateIndex\in\NonnegativeIntegers$.
Define
\[
 \begin{aligned}
 \NonInvariantSet _\CoordinateIndex &=\bigl\{\FiberPoint=\MeasuredSpace{\BaseCarrier }{\MetricSymbol }{\FirstMeasure }\in\MeasuredSpaces\bigm|
       \exists\GroupElement \in\Isom\MetricSpace{\BaseCarrier }{\MetricSymbol }\\
    &\qquad\qquad\ProkhorovDistance{\MetricSymbol }(\FirstMeasure ,\GroupElement \Pushforward\FirstMeasure )\geq2^{-\CoordinateIndex} \bigr\}.
 \end{aligned}
\]
We prove that
$\NonInvariantSet_\CoordinateIndex$
is closed.
Take a sequence
\[
 \FiberPoint_\SequenceIndex
 =\MeasuredSpace{\BaseCarrier_\SequenceIndex}{\MetricSymbol_\SequenceIndex}{\FirstMeasure_\SequenceIndex}
 \in\NonInvariantSet_\CoordinateIndex
\]
converging to
$\FiberPoint=\MeasuredSpace{\BaseCarrier}{\MetricSymbol}{\FirstMeasure}\in\MeasuredSpaces$.
For each
$\SequenceIndex$,
choose
$\GroupElement_\SequenceIndex\in\Isom\MetricSpace{\BaseCarrier_\SequenceIndex}{\MetricSymbol_\SequenceIndex}$
such that
\[
 \ProkhorovDistance{\MetricSymbol_\SequenceIndex}
 (\FirstMeasure_\SequenceIndex,
  (\GroupElement_\SequenceIndex)\Pushforward\FirstMeasure_\SequenceIndex)
 \geq2^{-\CoordinateIndex}.
\]
By \Cref{lem:common-measured-embedding},
there exist a compact metric space
$\MetricSpace{\AmbientCarrier}{\AmbientMetric}$
and isometric embeddings
$\FirstEmbedding_\SequenceIndex\colon\MetricSpace{\BaseCarrier_\SequenceIndex}{\MetricSymbol_\SequenceIndex}\to\MetricSpace{\AmbientCarrier}{\AmbientMetric}$
and
$\FirstEmbedding\colon\MetricSpace{\BaseCarrier}{\MetricSymbol}\to\MetricSpace{\AmbientCarrier}{\AmbientMetric}$
such that
\[
 \HausdorffDistance{\AmbientMetric}
 (\FirstEmbedding_\SequenceIndex(\BaseCarrier_\SequenceIndex),\FirstEmbedding(\BaseCarrier))\to0,
\]
and
\[
 (\FirstEmbedding_\SequenceIndex)\Pushforward\FirstMeasure_\SequenceIndex
 \to\FirstEmbedding\Pushforward\FirstMeasure
 \quad\text{weakly in }\ProbabilityMeasures(\AmbientCarrier).
\]
Identify the carriers with these images and the measures with these pushforwards.
For each
$\SequenceIndex$,
on the embedded carrier we use
$\GroupElement_\SequenceIndex$
to denote the conjugate
$\FirstEmbedding_\SequenceIndex\circ\GroupElement_\SequenceIndex\circ\FirstEmbedding_\SequenceIndex^{-1}$.
By \Cref{lem:isometry-limits},
we obtain a strictly increasing sequence of indices
$n_\ell$
and an isometry
$\GroupElement\in\Isom\MetricSpace{\BaseCarrier}{\MetricSymbol}$
for which
\[
 (\GroupElement_{n_\ell})\Pushforward\FirstMeasure_{n_\ell}
 \to\GroupElement\Pushforward\FirstMeasure
 \quad\text{weakly}.
\]
By \Cref{lem:prokhorov-isometric-embedding},
isometric embeddings preserve the L\'evy--Prokhorov distance.
Thus the triangle inequality on
$\ProbabilityMeasures(\AmbientCarrier)$
yields
\[
 \begin{aligned}
 &\left|
 \ProkhorovDistance{\MetricSymbol_{n_\ell}}
 (\FirstMeasure_{n_\ell},(\GroupElement_{n_\ell})\Pushforward\FirstMeasure_{n_\ell})
 -\ProkhorovDistance{\MetricSymbol}(\FirstMeasure,\GroupElement\Pushforward\FirstMeasure)
 \right|\\
 &\quad=\left|
 \ProkhorovDistance{\AmbientMetric}
 (\FirstMeasure_{n_\ell},(\GroupElement_{n_\ell})\Pushforward\FirstMeasure_{n_\ell})
 -\ProkhorovDistance{\AmbientMetric}(\FirstMeasure,\GroupElement\Pushforward\FirstMeasure)
 \right|\\
 &\quad\leq
 \ProkhorovDistance{\AmbientMetric}(\FirstMeasure_{n_\ell},\FirstMeasure)
 +\ProkhorovDistance{\AmbientMetric}
 ((\GroupElement_{n_\ell})\Pushforward\FirstMeasure_{n_\ell},\GroupElement\Pushforward\FirstMeasure)
 \longrightarrow0.
 \end{aligned}
\]
Both terms on the right tend to
$0$
because
$\FirstMeasure_{n_\ell}\to\FirstMeasure$
and
$(\GroupElement_{n_\ell})\Pushforward\FirstMeasure_{n_\ell}\to\GroupElement\Pushforward\FirstMeasure$
weakly in
$\ProbabilityMeasures(\AmbientCarrier)$,
and
$\ProkhorovDistance{\AmbientMetric}$
induces weak convergence.
Consequently,
\[
 \ProkhorovDistance{\MetricSymbol }(\FirstMeasure ,\GroupElement \Pushforward\FirstMeasure )
 =\lim_{\ell\to\infty} \ProkhorovDistance{\MetricSymbol _{n_\ell}}(\FirstMeasure _{n_\ell} ,(\GroupElement _{n_\ell} )\Pushforward\FirstMeasure _{n_\ell} )\geq2^{-\CoordinateIndex} .
\]
Thus
$\FiberPoint\in\NonInvariantSet_\CoordinateIndex$,
and hence
$\NonInvariantSet_\CoordinateIndex$
is closed.
Since the L\'evy--Prokhorov distance is a metric,
we see that
\[
 \mathscr S_{\mathrm{inv}}
 =\bigcap_{\CoordinateIndex\in\NonnegativeIntegers}
       (\MeasuredSpaces\setminus\NonInvariantSet_\CoordinateIndex).
\]
Each set in this intersection is open.
\end{proof}

\begin{theorem}\label{prop:invariant-measured-polish}
\paraid{p-a10d59349312}
The space
$\InvariantMeasuredSpaces$
is Polish.
\end{theorem}

\begin{proof}
\paraid{p-463534d4ce84}
By \Cref{lem:full-support-gdelta,lem:invariant-measures-gdelta},
\[
 \InvariantMeasuredSpaces
 =\mathscr S_{\mathrm{full}}\cap\mathscr S_{\mathrm{inv}}
\]
is a
$\CountableIntersectionType$
subset of
$\MeasuredSpaces$.
\Cref{thm:ghp-polish} states that
$\MeasuredSpaces$
is Polish,
so this subspace is Polish.
\end{proof}

\paraid{p-47cbaeca7f5a}
Define
\[
 \MeasureProjection\colon\InvariantMeasuredSpaces\to\GHSpace
\]
by
\[
 \MeasureProjection\MeasuredSpace{\BaseCarrier}{\MetricSymbol}{\FirstMeasure}
 =\MetricSpace{\BaseCarrier}{\MetricSymbol}.
\]
Namely, this map forgets the measure.
\begin{lemma}\label{lem:invariant-full-support}
\paraid{p-ab6dab25fbcd}
Every nonempty compact metric space
$\MetricSpace{\BaseCarrier}{\MetricSymbol}$
admits an invariant probability with full support.
\end{lemma}

\begin{proof}
\paraid{p-fee9db58d6fe}
Fix a dense sequence
$\{\BasePoint_\CoordinateIndex\}_{\CoordinateIndex\in\NonnegativeIntegers}$
in
$\BaseCarrier$
and put
\[
 \ApproximatingProbability
 =\sum_{\CoordinateIndex\in\NonnegativeIntegers}
   2^{-(\CoordinateIndex+1)}\DiracProbability_{\BasePoint_\CoordinateIndex}.
\]
Since
$\sum_{\CoordinateIndex\in\NonnegativeIntegers}2^{-(\CoordinateIndex+1)}=1$,
this is a probability on
$\BaseCarrier$.
By the density of
$\{\BasePoint_\CoordinateIndex\}_{\CoordinateIndex\in\NonnegativeIntegers}$,
we see that
$\supp(\ApproximatingProbability)=\BaseCarrier$.

\paraid{p-b8095b9a1aa1}
Put
$\IsometryGroup=\Isom\MetricSpace{\BaseCarrier}{\MetricSymbol}$.
By \Cref{lem:compact-isometry-group},
the group
$\IsometryGroup$
is compact and metrizable.
By the Haar theorem (\Cref{thm:haar}),
we obtain a normalized left invariant probability
$\HaarProbability _\IsometryGroup $
on
$\IsometryGroup $.
In the iterated integral below,
the outer variable
$\GroupElement\in\IsometryGroup$
is integrated against
$\HaarProbability_\IsometryGroup$,
and the inner variable
$\BasePoint\in\BaseCarrier$
is integrated against
$\ApproximatingProbability$.
Define the average
$\overline\ApproximatingProbability$
by requiring that for every
$\TestFunction\in\ContinuousFunctions(\BaseCarrier)$,
\[
 \int_\BaseCarrier \TestFunction \,\IntegrationDifferential \overline\ApproximatingProbability
 =\int_\IsometryGroup \int_\BaseCarrier \TestFunction (\GroupElement \BasePoint )\,\IntegrationDifferential \ApproximatingProbability (\BasePoint )\,\IntegrationDifferential \HaarProbability _\IsometryGroup (\GroupElement ).
\]
By the Riesz--Markov--Kakutani theorem (\Cref{thm:riesz-markov-kakutani}),
this positive normalized linear functional represents a unique probability on
$\BaseCarrier $.
Left invariance of
$\HaarProbability _\IsometryGroup $
implies
$\TranslatingIsometry \Pushforward\overline\ApproximatingProbability =\overline\ApproximatingProbability $
for every
$\TranslatingIsometry \in \IsometryGroup $.
For every nonempty open set
$\OpenSubset\subset\BaseCarrier$,
choose a nonzero continuous function
$\TestFunction\colon\BaseCarrier\to[0,1]$
that vanishes outside
$\OpenSubset$.
For each
$\GroupElement\in\IsometryGroup$,
the function
$\TestFunction\circ\GroupElement$
is nonzero and nonnegative,
so its integral against the full-support measure
$\ApproximatingProbability$
is positive.
The defining identity for the average therefore implies
\[
 \overline\ApproximatingProbability(\OpenSubset)
 \geq\int_\IsometryGroup\int_\BaseCarrier
       \TestFunction(\GroupElement\BasePoint)
       \,\IntegrationDifferential\ApproximatingProbability(\BasePoint)
       \,\IntegrationDifferential\HaarProbability_\IsometryGroup(\GroupElement)>0.
\]
Hence
$\MeasuredSpace{\BaseCarrier }{\MetricSymbol }{\overline\ApproximatingProbability }\in\InvariantMeasuredSpaces$,
as claimed.
\end{proof}

\begin{proposition}\label{lem:convergent-invariant-lifts}
\paraid{p-a0c55c41bf20}
Let
$\FiberPoint=\MeasuredSpace{\BaseCarrier}{\MetricSymbol}{\FirstMeasure}\in\InvariantMeasuredSpaces$
and assume that
$\MetricSpace{\BaseCarrier_\SequenceIndex}{\MetricSymbol_\SequenceIndex}
 \to\MetricSpace{\BaseCarrier}{\MetricSymbol}$
in
$\GHSpace$.
Then there are invariant probabilities
$\SecondMeasure_\SequenceIndex$
with full support on
$\BaseCarrier_\SequenceIndex$
such that
\[
 \MeasuredSpace{\BaseCarrier_\SequenceIndex}{\MetricSymbol_\SequenceIndex}{\SecondMeasure_\SequenceIndex}
 \longrightarrow\FiberPoint
 \quad\text{in }\InvariantMeasuredSpaces.
\]
\end{proposition}

\begin{proof}
\paraid{p-c601850135cd}
By \Cref{lem:common-gh-embedding},
choose isometric embeddings into a compact
$\MetricSpace{\AmbientCarrier }{\AmbientMetric }$
and identify the spaces with their images.
Let $\iota_\SequenceIndex\colon\BaseCarrier_\SequenceIndex\hookrightarrow\AmbientCarrier$
and $\iota\colon\BaseCarrier\hookrightarrow\AmbientCarrier$ be the inclusions.
By the Hausdorff continuity of the probability-measure functor
\cite[arXiv v5, Lemma~2.13 and Example~2.1(iii)]{Khezeli2023Framework},
there are probabilities
$\FirstMeasure_\SequenceIndex\in\ProbabilityMeasures(\BaseCarrier_\SequenceIndex)$
such that
$(\iota_\SequenceIndex)_*\FirstMeasure_\SequenceIndex\to\iota_*\FirstMeasure$
weakly on $\AmbientCarrier$.
Define
$\bar\mu_\SequenceIndex=(\iota_\SequenceIndex)_*\FirstMeasure_\SequenceIndex$
and $\bar\mu=\iota_*\FirstMeasure$ in $\ProbabilityMeasures(\AmbientCarrier)$.
Then
\[
 \ProkhorovDistance{\AmbientMetric}(\bar\mu_\SequenceIndex,\bar\mu)\longrightarrow0.
\]
Thus $\bar\mu_\SequenceIndex\to\bar\mu$ weakly in
$\ProbabilityMeasures(\AmbientCarrier)$.
Choose a probability
$\FullSupportProbability_\SequenceIndex$
with full support on each
$\BaseCarrier_\SequenceIndex$,
for example by assigning positive weights summing to one to a dense sequence.
For
$0<\MixingWeight_\SequenceIndex<1$
with
$\MixingWeight_\SequenceIndex\to0$,
put
\[
 \ApproximatingProbability_\SequenceIndex
 =(1-\MixingWeight_\SequenceIndex)\FirstMeasure_\SequenceIndex
   +\MixingWeight_\SequenceIndex\FullSupportProbability_\SequenceIndex.
\]
This probability has full support.
For every real continuous function
$\TestFunction$
on
$\AmbientCarrier$,
\[
 \left|\int_{\BaseCarrier_\SequenceIndex}(\TestFunction\circ\iota_\SequenceIndex)\,\IntegrationDifferential\ApproximatingProbability_\SequenceIndex
       -\int_{\BaseCarrier_\SequenceIndex}(\TestFunction\circ\iota_\SequenceIndex)\,\IntegrationDifferential\FirstMeasure_\SequenceIndex\right|
 \leq2\MixingWeight_\SequenceIndex\norm{\TestFunction}_\infty\longrightarrow0.
\]
Hence
$(\iota_\SequenceIndex)_*\ApproximatingProbability_\SequenceIndex\to\bar\mu$
weakly.

\paraid{p-0cb8c42117db}
Average
$\ApproximatingProbability _\SequenceIndex $
with respect to the normalized Haar probability on
$\IsometryGroup _\SequenceIndex =\Isom\MetricSpace{\BaseCarrier _\SequenceIndex }{\MetricSymbol _\SequenceIndex}$
and denote the resulting probability by
$\SecondMeasure _\SequenceIndex $.
The averaging construction in \Cref{lem:invariant-full-support}
shows that these probabilities are invariant and have full support.
For each
$\TestFunction \in\ContinuousFunctions(\AmbientCarrier )$,
we claim that
\begin{equation}\label{eq:uniform-orbit-measures}
 \sup_{\GroupElement \in \IsometryGroup _\SequenceIndex }
 \left|\int_{\BaseCarrier_\SequenceIndex}(\TestFunction\circ\iota_\SequenceIndex) \,\IntegrationDifferential (\GroupElement \Pushforward\ApproximatingProbability _\SequenceIndex )-\int_\BaseCarrier(\TestFunction\circ\iota) \,\IntegrationDifferential \FirstMeasure \right|
 \longrightarrow0.
\end{equation}
For the sake of contradiction,
suppose that \eqref{eq:uniform-orbit-measures} fails.
Choose
$\ErrorTolerance>0$
and a strictly increasing sequence of indices
$\{n_k\}_{k\in\NonnegativeIntegers}$
and for each
$k\in\NonnegativeIntegers$
choose
$\GroupElement_{n_k}\in\IsometryGroup_{n_k}$
satisfying
\begin{equation}\label{eq:orbit-integral-lower-bound}
 \left|\int_{\BaseCarrier_{n_k}}(\TestFunction\circ\iota_{n_k})\,
       \IntegrationDifferential((\GroupElement_{n_k})\Pushforward\ApproximatingProbability_{n_k})
       -\int_\BaseCarrier(\TestFunction\circ\iota)\,\IntegrationDifferential\FirstMeasure\right|
 \geq\ErrorTolerance.
\end{equation}
By \Cref{lem:isometry-limits},
there exist a strictly increasing sequence of indices
$\{k(j)\}_{j\in\NonnegativeIntegers}$
and an isometry
$\GroupElement\in\Isom\MetricSpace{\BaseCarrier}{\MetricSymbol}$
such that
\[
 (\iota_{n_{k(j)}}\circ\GroupElement_{n_{k(j)}})_*\ApproximatingProbability_{n_{k(j)}}
 \to(\iota\circ\GroupElement)_*\FirstMeasure=\bar\mu.
\]
Apply this weak convergence to the fixed function
$\TestFunction\in\ContinuousFunctions(\AmbientCarrier)$.
By the pushforward integral identity,
\[
 \left|\int_{\BaseCarrier_{n_{k(j)}}}(\TestFunction\circ\iota_{n_{k(j)}})
 \,\IntegrationDifferential((\GroupElement_{n_{k(j)}})_*\ApproximatingProbability_{n_{k(j)}})
 -\int_\BaseCarrier(\TestFunction\circ\iota)\,\IntegrationDifferential\FirstMeasure\right|
 \longrightarrow0.
\]
This contradicts the lower bound $\ErrorTolerance>0$ in
\eqref{eq:orbit-integral-lower-bound}.
Therefore,
averaging \eqref{eq:uniform-orbit-measures} over
$\IsometryGroup _\SequenceIndex $
yields
$(\iota_\SequenceIndex)_*\SecondMeasure_\SequenceIndex\to\bar\mu$
weakly.
Define
\[
 \FiberPoint_\SequenceIndex
 :=\MeasuredSpace{\BaseCarrier_\SequenceIndex}{\MetricSymbol_\SequenceIndex}{\SecondMeasure_\SequenceIndex}.
\]
Consequently,
\[
 \FiberPoint_\SequenceIndex\longrightarrow\FiberPoint,
\]
and
\[
 \MeasureProjection(\FiberPoint _\SequenceIndex )=\MetricSpace{\BaseCarrier _\SequenceIndex }{\MetricSymbol _\SequenceIndex}.
\]
Thus
$\{\FiberPoint_\SequenceIndex\}_{\SequenceIndex\in\NonnegativeIntegers}$
is the required sequence of invariant lifts.
This finishes the proof.
\end{proof}

\begin{theorem}\label{prop:open-invariant-projection}
\paraid{p-909c842b7494}
The map
\[
 \MeasureProjection\colon\InvariantMeasuredSpaces\to\GHSpace
\]
defined by
\[
 \MeasureProjection\MeasuredSpace{\BaseCarrier}{\MetricSymbol}{\FirstMeasure}
 =\MetricSpace{\BaseCarrier}{\MetricSymbol},
\]
is a continuous open surjection.
\end{theorem}

\begin{proof}
\paraid{p-34f1f298a7ea}
The inequality
\[
 \GHDistance(\MetricSpace{\BaseCarrier }{\MetricSymbol },\MetricSpace{\ComparisonCarrier }{\ComparisonMetric })
 \leq\GHPDistance(\MeasuredSpace{\BaseCarrier }{\MetricSymbol }{\FirstMeasure },\MeasuredSpace{\ComparisonCarrier }{\ComparisonMetric }{\SecondMeasure })
\]
shows that
$\MeasureProjection$
is continuous.
By \Cref{lem:invariant-full-support},
the map
$\MeasureProjection$
is surjective.
To prove openness,
let
$\OpenSubset \subset\InvariantMeasuredSpaces$
be open,
and let
$\FiberPoint \in \OpenSubset $.
For the sake of contradiction,
suppose that
$\MeasureProjection(\OpenSubset )$
contains no neighborhood of
$\MeasureProjection(\FiberPoint )$.
For each
$\SequenceIndex\in\NonnegativeIntegers$,
choose
\[
 \MetricSpace{\BaseCarrier _\SequenceIndex }{\MetricSymbol _\SequenceIndex}\in
 \GHball(\MeasureProjection(\FiberPoint ),2^{-\SequenceIndex} )\setminus\MeasureProjection(\OpenSubset ).
\]
By \Cref{lem:convergent-invariant-lifts},
there exist
$\FiberPoint_\SequenceIndex\in\InvariantMeasuredSpaces$
such that
\[
 \MeasureProjection(\FiberPoint_\SequenceIndex)
 =\MetricSpace{\BaseCarrier_\SequenceIndex}{\MetricSymbol_\SequenceIndex},
\]
and
\[
 \FiberPoint_\SequenceIndex\to\FiberPoint.
\]
Since
$\OpenSubset$
is an open neighborhood of
$\FiberPoint$,
we obtain
$\FiberPoint _\SequenceIndex \in \OpenSubset $
for all sufficiently large
$\SequenceIndex $.
This implies
$\MetricSpace{\BaseCarrier _\SequenceIndex }{\MetricSymbol _\SequenceIndex}\in\MeasureProjection(\OpenSubset )$,
a contradiction.
Thus
$\MeasureProjection(\OpenSubset )$
is open.
\end{proof}
